\documentclass[12pt]{article}

%Dois condutores infinitos
%Faro, Mon Mar  3 21:36:38 WET 2014
%Written by Tordar

\usepackage{graphicx}
\usepackage{pst-all}



\begin{document}

\pagestyle{empty}

\begin{pspicture}(-5,-5)(5,5)
\psline[linewidth=1.5mm,linecolor=black](-4,-3)(-4,3)
\psline[linewidth=1.5mm,linecolor=black]( 4,-3)( 4,3)
\psline[linewidth=1.5mm,linecolor=black,linestyle=dashed](-4,-5)(-4,-3)
\psline[linewidth=1.5mm,linecolor=black,linestyle=dashed](-4, 3)(-4, 5)
\psline[linewidth=1.5mm,linecolor=black,linestyle=dashed]( 4,-5)( 4,-3)
\psline[linewidth=1.5mm,linecolor=black,linestyle=dashed]( 4, 3)( 4, 5)
\psline[linewidth=0.7mm,linecolor=black,linestyle=dashed](0,-5)(0,5)
\psline[linewidth=1.5mm,linecolor=blue,arrows=->](-5,-2)(-5,2)
\psline[linewidth=1.5mm,linecolor=blue,arrows=->]( 5,-2)( 5,2)
\psline[linewidth=1.5mm,linecolor=blue,arrows=<->](-4,0)(0,0)
\psline[linewidth=1.5mm,linecolor=blue,arrows=<->](0,0)(4,0)
\rput{0}(-5.5,0){\scalebox{1.7}{$I$}}
\rput{0}( 5.5,0){\scalebox{1.7}{$I$}}
\rput{0}(-2,0.4){\scalebox{1.5}{$a$}}
\rput{0}( 2,0.4){\scalebox{1.5}{$a$}}
\end{pspicture}

\end{document}
